\section{All late corrections of a factorial Lambert inverse}
\label{llc:section}

The preceding Lambert chart resolves a logarithmic marginal scale. A
different Lambert structure arises when the forward tail itself has
factorial growth. Here the small-index coefficients of that tail determine
every inverse-index correction by a finite polynomial calculation.
The leading hyperbolic amplitudes and their first correction were proved
in the reversion continuation \cite{llc:reversion}, resolving the
conjectures of \cite{llc:inverse}. We give a uniform extension to every
fixed order, motivated by its all-order question. The experimental
calculation of the second correction guides the finite formula; the
factorial estimates below justify the full infinite asymptotic statement.

\subsection{The inverse equation and finite coefficient extraction}

Put $A=4\pi^2$ and let $d(n)$ count positive divisors. Mellin inversion
of the double Lambert sum gives
\[
 \mathcal F(t)=\sum_{n\ge1}\frac{d(n)}n e^{-nt},\qquad
 \Gamma(w)\zeta(1+w)^2t^{-w}
\]
as function and Mellin integrand. With the Stieltjes convention
$\zeta(1+w)=w^{-1}+\sum_{j\ge0}(-1)^j\gamma_jw^j/j!$, its
triple pole at zero contributes
\[
 \frac12\bigl(\log(1/t)+\gamma\bigr)^2+C_* ,\qquad
 C_* =\frac{\zeta(2)}2-\gamma^2-2\gamma_1.
\]
The remaining residues define the formal tail
\begin{equation}\label{llc:tail}
 R(q)=\sum_{m\ge1}r_mq^m,\quad c=r_1=-\frac14,\quad
 r_{2m}=\frac{4\Gamma(2m)\zeta(2m)^2}{2mA^{2m}},\quad
 r_{2m+1}=0\quad(m\ge1).
\end{equation}
Equivalently $r_m=(-1)^m\zeta(1-m)^2/m!$, so all coefficients
are rational; in particular $r_2=1/288$. The exact real completion
is treated in \cite{llc:inverse}; the coefficient argument here uses
only \eqref{llc:tail}.

For a large target $Y$, set $X=\sqrt{2(Y-C_*)}$, $u=X^{-1}$,
$q=e^{\gamma-X}$ and $t=qe^{-V}$. The corresponding formal equation is
\begin{equation}\label{llc:inverse-equation}
 V+\frac u2V^2+uR(qe^{-V})=0,\qquad
 V(q,u)=\sum_{n\ge1}C_n(u)q^n.
\end{equation}
In the following theorem $q$ and $u$ are independent variables.
Formal implicit inversion gives a unique solution over $\mathbb Q[u]$.
Introduce the finite polynomial
\begin{equation}\label{llc:H}
 H(w,k,u)=\sum_{v=0}^{k-1}\binom{k-1}{v}(k)_v
                       (u/2)^v w^v,\qquad k\in\mathbb Z_{>0},
\end{equation}
where $(k)_v$ is the rising factorial. Auxiliary-variable Lagrange
inversion yields the exact identity
\begin{equation}\label{llc:finite}
 C_n(u)=-\sum_{k=1}^n\frac{u^k}{k!}n^{k-1}
                     [q^n]R(q)^k H(1/n,k,u).
\end{equation}
Indeed, apply Lagrange inversion to
$V=-z uR(qe^{-V})(1+uV/2)^{-1}$. The coefficient of $z^kq^n$
is $(-u)^k[q^n]R^k/k$ times
$[v^{k-1}]e^{-nv}(1+uv/2)^{-k}$, which is the displayed term.
Only $k\le n$ occurs, so setting $z=1$ is a formal operation.

\subsection{A finite polynomial for every correction}

Let $S_r(L)$ be the Faulhaber polynomial equal to
$\sum_{i=1}^L i^r$ for nonnegative integers $L$. Define the formal series
\begin{equation}\label{llc:B}
 B(w,L)=\exp\left(\sum_{r\ge1}\frac{L^r+S_r(L)}r w^r\right).
\end{equation}
For integer $L\ge0$ this is
$(1-Lw)^{-1}\prod_{i=1}^L(1-iw)^{-1}$.
When extracting a fixed power of $w$, we interpret \eqref{llc:H}
polynomially in $k$, using falling factorials for $\binom{k-1}{v}$.
Thus both expressions make sense as finite coefficient operations
even before an integer value is substituted.

Define the linear Poisson polynomial transform by
\begin{equation}\label{llc:poisson}
 \mathcal P_z(L^r)=e^{-z}\sum_{\ell\ge0}\frac{\ell^r z^\ell}{\ell!}
                 =\sum_{a=0}^r {r\brace a}z^a.
\end{equation}
For $r=0$ its value is $1$, with $0^0=1$ in the sum.
The Stirling-number formula makes the operation finite.

Fix $j\ge0$. Range over nonnegative multiplicities
$\mathbf a=(a_2,\ldots,a_{j+1})$ satisfying
\[
 h=\sum a_m,\qquad s=\sum m a_m,\qquad
 d=s-h=\sum(m-1)a_m\le j.
\]
The empty tuple is allowed. For $\epsilon\in\{1,-1\}$ set
$\theta=Au/4$ and $z_\epsilon=\epsilon Acu=-\epsilon\theta$.
The correction polynomial is
\begin{equation}\label{llc:T}
 \begin{split}
 T_j^\epsilon(u)={}&\sum_{\mathbf a:\ d\le j}
 u^h(\epsilon A)^s\prod_{m=2}^{j+1}\frac{r_m^{a_m}}{a_m!}\,
 \mathcal P_{z_\epsilon}\!\left(
 [w^{j-d}]B(w,L+s)H(w,L+h+1,u)\right).
 \end{split}
\end{equation}
The transform acts on $L$ only. Every sum and coefficient extraction
in \eqref{llc:T} is finite; its only forward coefficients are
$r_1,\ldots,r_{j+1}$. In particular $T_0^\epsilon=1$.

\begin{theorem}[All fixed-order Lambert corrections]\label{llc:all-orders}
For each integer $J\ge0$ and $U>0$, uniformly for complex $|u|\le U$,
\begin{equation}\label{llc:expansion}
 \begin{split}
 C_n(u)={}&-\frac{2u\Gamma(n)}{nA^n}
 \sum_{\epsilon=\pm1}\epsilon^n e^{-\epsilon\theta}
                  \sum_{j=0}^J\frac{T_j^\epsilon(u)}{n^j}
 +O_{J,U}\!\left(\frac{|u|\Gamma(n)}{A^n n^{J+2}}\right).
 \end{split}
\end{equation}
This is an additive expansion, including parameter values where a
parity amplitude vanishes. It asserts an asymptotic expansion at every
fixed order, with no assertion of convergence as $J\to\infty$.
\end{theorem}

\subsection{Uniform control of the omitted configurations}

We supply the estimates that allow the finite calculation to be summed
over all factor counts in \eqref{llc:finite}.

\begin{lemma}[Factorial vertices]\label{llc:vertices}
For positive integers $N\ge p$,
\begin{equation}\label{llc:convolution}
 \sum_{\substack{\ell_1+\cdots+\ell_p=N\\\ell_i\ge1}}
                 \prod_i(\ell_i-1)!\le3^{p-1}(N-p)!.
\end{equation}
Fix $J\ge0$. If the configurations with
$N-\max_i\ell_i\le J$ are removed, then the remaining sum is at most
$D_J^p(N-J-2)!$, for all sufficiently large $N$ depending only on $J$,
with a constant $D_J\ge1$ independent of $p$ and $N$.
\end{lemma}
\begin{proof}
The two-factor inequality
$\sum_{a=0}^M a!(M-a)!\le3M!$ follows by dividing by $M!$:
the endpoint terms sum to $2$ and each interior reciprocal binomial
coefficient is at most $1/M$. The cases $M=0,1$ are immediate.
Induction proves \eqref{llc:convolution}.

If $p\ge J+2$, this bound already suffices since
$(N-p)!\le(N-J-2)!$. If $p=1$ nothing remains.
For $2\le p\le J+1$, choose a maximal coordinate and put
$d=N-\max_i\ell_i$. The remainder is at most
\[
 p3^{p-2}\sum_{d=J+1}^{\lfloor N-N/p\rfloor}
                    (d-p+1)!(N-d-1)!.
\]
Write $M=N-p$ and $b=J+2-p\ge1$. For sufficiently large $N$,
both factorial indices in this sum lie between $b$ and $M-b$.
For each fixed $b$,
\begin{equation}\label{llc:factorial-tail}
 \sum_{t=b}^{M-b}t!(M-t)!\le K_b(M-b)!.
\end{equation}
To verify this uniformly, in the lower outer third the ratio of
successive terms is $(t+1)/(M-t)\le3/4$ for large $M$.
Sum that geometric majorant starting at $b!(M-b)!$ and use symmetry
for the upper third. The central terms are bounded by
$b!(M-b)!(3/4)^{\lfloor M/3\rfloor-b}$, times a fixed factor;
their number is at most $M+1$. This is uniformly bounded after division
by $(M-b)!$. The finitely many small $M$ can be absorbed in $K_b$.
Equation \eqref{llc:factorial-tail} now gives $(N-J-2)!$.
Only finitely many $p\le J+1$ remain, so a common $D_J$ suffices.
\end{proof}

\begin{proof}[Proof of Theorem~\ref{llc:all-orders}]
Choose $K\ge1$ such that $A|r_1|\le K$ and
$A^m|r_m|\le K(m-2)!$ for $m\ge2$. Such a $K$ exists by
\eqref{llc:tail}. For a composition $m_1+\cdots+m_k=n$, put
$\ell_i=m_i-1$ and $N=n-k$. Choosing the $p$ positive increments
and applying \eqref{llc:convolution} gives, for $k<n$,
\begin{equation}\label{llc:global-bound}
 A^n\bigl|[q^n]R^k\bigr|\le(4K)^k\Gamma(n-k).
\end{equation}
Absolute values may also be placed inside the coefficient sum.
Since $(k)_v\le(2k)^v$ for $v\le k-1$,
$|H(1/n,k,u)|\le(1+U)^k$ for $k\le n$.
After division by $\Gamma(n)/(nA^n)$, the terms with
$k\le n/2$ are bounded by $D^k/k!$ for a constant $D=D_U$:
indeed $n^k\Gamma(n-k)/\Gamma(n)\le2^k$.
The sum over $n/2<k<n$ is at most
\[
 n\max(1,2eD)^n\frac{\Gamma(\lceil n/2\rceil)}{\Gamma(n)},
\]
using $n^k/k!\le(2e)^k$. It is smaller than every inverse power
of $n$. The $k=n$ term, where all indices are $1$, has the same
property directly from $n^n/(n!\Gamma(n))$. All these estimates
retain an overall factor $|u|$, by bounding $|u|^k$ by
$|u|\max(1,U)^{k-1}$.

For $k\le n/2$, retain exactly the configurations with one large
increment and total other increment at most $J$. Lemma~\ref{llc:vertices},
including the choices of positive positions, bounds the omitted
absolute coefficient sum by
$A^{-n}D_{J}^k\Gamma(n-k-J-1)$, with an enlarged constant depending
also on $K$. For sufficiently large $n$,
\[
 \frac{n^k\Gamma(n-k-J-1)}{\Gamma(n)}
 \le n^{-J-1}3^{k+J+1}.
\]
Together with the $k!$ denominator and the bound on $H$, this proves
that the total omitted contribution is
$O_{J,U}(|u|\Gamma(n)/(A^n n^{J+2}))$.

In a retained configuration let $\ell$ be the number of indices $1$,
and let $a_m$ count the other small indices. Then
$k=\ell+h+1$ and the large index is $m=n-\ell-s$.
The large increment is unique once $n$ is sufficiently large, since
$N>2J$. The contribution to $[q^n]R^k$ is exactly
\[
 \frac{k!}{\ell!\prod a_m!}\,
 c^\ell\prod r_m^{a_m}\,r_{n-\ell-s}.
\]
Here $n-\ell-s\ge n/2+1-J$. Replacing $\zeta(n-\ell-s)^2$
by $1$ in the large even coefficient has an exponentially small
relative error, uniformly in these configurations. The odd large
coefficients vanish exactly. Thus use
\[
 r_m=\frac{2\Gamma(m)}{mA^m}\sum_{\epsilon=\pm1}\epsilon^m
\]
at a total exponentially small cost under the same factorial majorant.
The exact gamma factor is
\begin{equation}\label{llc:gamma-product}
 \frac{n^{L+1}\Gamma(n-L)}{(n-L)\Gamma(n)}=B(1/n,L),
 \qquad L=\ell+s.
\end{equation}
Consequently, for each small multiset the normalized action contribution
is the finite sum of
\[
 \frac{z_\epsilon^\ell}{\ell!}\,
 u^h(\epsilon A)^s\prod\frac{r_m^{a_m}}{a_m!}\,
 n^{-d}B(1/n,\ell+s)H(1/n,\ell+h+1,u).
\]
Its common prefactor is $-2u\epsilon^n\Gamma(n)/(nA^n)$.

Expand the product $BH$ through degree $J-d$ in $1/n$.
Here is a uniform remainder estimate for this step. We have
$L\le n/2+J$, so $L/n\le3/4$ for large $n$.
On $0\le w\le1/n$, the product defining $B$ is at most $4^{L+1}$.
The derivatives of its logarithm of order $a\ge1$ are bounded by
a constant times $(1+L)^{a+1}$, with constants depending on $a$.
Its fixed-order Taylor remainder is therefore bounded by
$n^{-b-1}$ times a polynomial in $L$ times $4^{L+1}$ when truncated
at order $b$. For $H$, the terms beyond degree $b$ are bounded by
\[
 \sum_{v=b+1}^{k-1}\binom{k-1}{v}(kU/n)^v
 \le (k/n)^{b+1}(1+U)^k.
\]
The fixed retained coefficients of $H$ are polynomials in $k,u$.
These bounds give the same polynomial-times-exponential remainder
for the product. It is summable against $|z_\epsilon|^\ell/\ell!$,
uniformly on $|u|\le U$. Extending the resulting finite $\ell$ sums
to infinity costs less than every inverse power of $n$, since only
fixed-degree polynomials in $\ell$ remain. The coefficient of $n^{-j}$
is now precisely the Poisson transform in \eqref{llc:T}, multiplied
by $e^{z_\epsilon}$. This proves \eqref{llc:expansion}.
\end{proof}

\subsection{The explicit second correction}

Let $b=uA^2/288$ and define
\begin{align}
 P_1&=\frac{1+u}{2}\theta^2+b,& D_1&=u+2,\label{llc:first}\\
 P_2&=\frac{(1+u)^2}{8}\theta^4
 +\frac{3u^2+11u+11+(u+1)b}{2}\theta^2
 +(u+5)b+\frac{b^2}{2},\label{llc:second-P}\\
 D_2&=\frac{6u^2+15u+10}{6}\theta^2+(2u+4)b+2u+3.
 \label{llc:second-D}
\end{align}
Then $T_1^\epsilon=P_1-\epsilon\theta D_1$ and
$T_2^\epsilon=P_2-\epsilon\theta D_2$.
For $E_n(\theta)=\cosh\theta$ at even $n$ and
$E_n(\theta)=-\sinh\theta$ at odd $n$, the expansion becomes
\begin{equation}\label{llc:second}
 \frac{nA^nC_n(u)}{4\Gamma(n)}
 =-u\left[E_n+\frac{P_1E_n+D_1\theta E_n'}n
              +\frac{P_2E_n+D_2\theta E_n'}{n^2}\right]
       +O_U(|u|n^{-3}).
\end{equation}
The primes differentiate $E_n$ with respect to its argument.

For an explicit check of the calculation, write
\[
 b_1=\frac{\ell(\ell+3)}2,\quad
 b_2=\frac{\ell(2\ell^2+9\ell+1)}{12}+\frac{b_1^2}{2},\quad
 h_1=\frac{u\ell(\ell+1)}2,\quad
 h_2=\frac{u^2\ell(\ell-1)(\ell+1)(\ell+2)}8.
\]
The configurations with only linear small indices give
$b_2+b_1h_1+h_2$ before the Poisson transform.
One small index $2$ gives
$b\{(\ell+2)(\ell+5)+u(\ell+2)(\ell+1)\}/2$;
two such indices give $b^2/2$. There is no small-index-$3$
term because $r_3=0$. Applying \eqref{llc:poisson} to these
expressions gives \eqref{llc:second-P}--\eqref{llc:second-D}.

\subsection{Higher polylogarithmic Borel boundary models}

Use the shifted Borel convention
$\widehat V_u(\xi)=\sum_{n\ge1}C_n(u)\xi^{n-1}/\Gamma(n)$.
For each fixed $J\ge0$, Theorem~\ref{llc:all-orders} gives the exact
decomposition in $|\xi|<A$
\begin{equation}\label{llc:borel}
 \widehat V_u(\xi)=-\frac{2u}{\xi}
 \sum_{\epsilon=\pm1}e^{-\epsilon\theta}
 \sum_{j=0}^J T_j^\epsilon(u)\Li_{j+1}(\epsilon\xi/A)
 +K_{J,u}(\xi).
\end{equation}
The apparent division by zero is removable. The remainder and its
first $J$ complex derivatives extend continuously to the closed disk,
locally uniformly in $u$: the $J$th differentiated remainder series
is bounded by a constant times $\sum n^{-2}$ there.
Thus \eqref{llc:second} supplies an explicit trilogarithmic model with
a twice continuously differentiable remainder. It extends the logarithmic
and dilogarithmic models of \cite{llc:reversion} to arbitrary finite
boundary regularity.

Since $\Li_1(z)=-\log(1-z)$, the two logarithmic coefficients at
$\xi=A$ and $\xi=-A$ are respectively
$2u e^{-\theta}/A$ and $-2u e^\theta/A$.
Both are nonzero for $u\ne0$. The continuous $J=0$ remainder therefore
proves genuine singularities at both endpoints and exact Borel radius
$A$; the ordinary series has radius zero. Finite-order boundary
regularity alone gives no continuation through those points and does
not determine the full local singularity on other sheets.

\subsection{Exact checks and numerical guidance}

The accompanying correction compiler constructs \eqref{llc:T} through
$j=4$ using rational arithmetic and Stirling polynomials. Independent
finite-product expansions check the Faulhaber coefficients at thirteen
integer values of $L$, and exact simplification checks both explicit
corrections. The comparison uses all forty fresh finite-sum coefficient
values from \cite{llc:reversion}, through $n=201$, including complex
coupling and a vanishing leading parity amplitude. For example, at
$u=0.1,n=200$, the absolute error in $nA^nC_n/(4\Gamma(n))$ decreases
from approximately $2.078\cdot10^{-3}$ at leading order to
$5.799\cdot10^{-5}$, $1.720\cdot10^{-6}$,
$5.779\cdot10^{-8}$ and $2.155\cdot10^{-9}$ through orders one to four.
These finite computations helped identify and check the correction
formula. They are numerical diagnostics, not interval certificates or
the proof of the uniform theorem.
